\par\Needspace{6\baselineskip}\medskip
\begin{center}\hypertarget{native-181-article135}{}\pdfbookmark[2]{135.}{article-135}135.\end{center}
Die vorangehenden Untersuchungen sollten die Grundsätze, auf denen unser Verfahren beruht, und
gewissermaßen seinen eigentlichen Kern deutlicher erkennen lassen. In der eigentlichen
Darstellung wird das Verfahren jedoch in einer ganz anderen Form erscheinen. Nach sehr häufiger
Anwendung können wir diese Form als die zweckmäßigste unter den verschiedenen von uns erprobten
Formen empfehlen. Da die Bestimmung einer unbekannten Bahn aus drei Beobachtungen immer auf
einige aufeinanderfolgende Hypothesen oder vielmehr Näherungen hinausläuft, ist es ein
beträchtlicher Gewinn, wenn man die Rechnung so einrichten kann, dass von Anfang an möglichst
viele Rechenschritte, die von $P$ und $Q$ unabhängig sind und allein von der Verbindung
bekannter Größen abhängen, von den Hypothesen selbst getrennt werden können. Diese
vorbereitenden Schritte, die allen Hypothesen gemeinsam sind, braucht man dann offenbar nur
einmal auszuführen, und bei den Hypothesen selbst fallen möglichst wenige Rechenschritte an.
Ebenso wichtig ist es, wenn man bei den einzelnen Hypothesen nicht bis zu den Elementen selbst
vordringen muss und ihre Berechnung der \hypertarget{native-182}{}letzten Hypothese vorbehalten kann. In beiderlei
Hinsicht scheint unser Verfahren, mit dessen Darstellung wir nun beginnen, nichts zu wünschen
übrig zu lassen.
